\clearpage
\fancyhead[L]{\small GERLING. SONNENFINSTERNIS.}
\hypertarget{native-627-lower}{}
\begin{center}\small Göttingische gelehrte Anzeigen. Stück 103. Seiten 1025..1030. 26. Juni 1820.\\[3pt]\rule{.72\linewidth}{.3pt}\end{center}

Göttingen. Zu der großen Sonnenfinsternis, die am 7. September dieses Jahres eintritt, hat Herr Professor \textsc{Gerling} in Marburg, der sie bereits vor acht Jahren in seiner Inauguraldissertation behandelt hatte (siehe unsere Anzeigen von 1812, S. 1889 [S. 542 dieses Bandes]), in einem Schreiben an Herrn Hofrat \textsc{Gauss} vom 23. Mai die Ergebnisse seiner wiederholten und weiter ausgeführten Untersuchungen mitgeteilt. Da diese Finsternis im größten Teil Deutschlands ringförmig erscheinen und für diese Gegenden auf lange Zeit das bemerkenswerteste Phänomen dieser Art sein wird, glauben wir, vielen Lesern dieser Blätter \SourceLeaf{628} durch die Veröffentlichung eines Auszugs aus jenen Mitteilungen einen willkommenen Dienst zu erweisen.

Die folgenden, den Sonnentafeln von \textsc{Carlini} und den Mondtafeln von \textsc{Bürg} entnommenen Sonnen- und Mondpositionen dienten als Grundlage für die Berechnung:

7. September 1820. Mittlere Pariser Zeit.

\begin{center}
\small\setlength{\tabcolsep}{2pt}\begin{tabular}{l|rrr|rrr|rrr}
 & \multicolumn{3}{c|}{\shortstack{Länge der\\Sonne.}} & \multicolumn{3}{c|}{\shortstack{Länge des\\Mondes.}} & \multicolumn{3}{c}{\shortstack{Nördl. Breite\\des Mondes.}} \\
\hline
$11^{\mathrm u}$ Vorm. & $164^\circ$ & $40'$ & $26''4$ & $163^\circ$ & $19'$ & $53''6$ & $0^\circ$ & $52'$ & $45''0$ \\
2 Nachm. & & $47$ & $44.0$ & $164$ & $48$ & $17.7$ & & $44$ & $41.6$ \\
5 & & $55$ & $1.7$ & $166$ & $16$ & $38.3$ & & $36$ & $34.9$ \\
\end{tabular}
\end{center}

\begin{center}\small\setlength{\tabcolsep}{3pt}
\begin{tabular}{lr}
Horizontalparallaxe des Mondes . . . . . . . . . . & $53'\ 55''5$ \\
\hspace{1em}---\hspace{1em}der Sonne . . . . . . . . . . & $8.7$ \\
Halbmesser des Mondes . . . . . . . . . . . . . & $14'\ 43.1$ \\
\hspace{1em}---\hspace{1em}der Sonne . . . . . . . . . . . . . & $15\ 54.8$ \\
Schiefe der Ekliptik . . . . . . . . . . . . . & $23^\circ\ 27\ 56.2$ \\
\end{tabular}
\begin{tabular}{ll}
Die Konjunktion erfolgt demnach um . . . . . . & $1^{\mathrm u}\ 58^{\mathrm m}\ 46^{\mathrm s}3$ mittl. Zeit in Paris \\
in der Länge . . . . . . . . . . . . . . . & $164^\circ\ 47'\ 41''$ \\
\end{tabular}
\end{center}

Für die folgenden einzelnen Orte hat Herr Professor \textsc{Gerling} zur Erleichterung der Beobachtungen die wichtigsten Zeitpunkte der Erscheinung im Voraus berechnet:

\begin{center}\small\setlength{\tabcolsep}{3pt}
\begin{tabular}{l|rrr|rrr|rrr}
 &  \multicolumn{3}{c|}{Anfang.}  &  \multicolumn{3}{c|}{Mitte.}  &  \multicolumn{3}{c|}{Ende.}  \\
\hline
Berlin  &  $1^{\mathrm u}$  &  $31^{\mathrm m}$  &  $55^{\mathrm s}$  &  $2^{\mathrm u}$  &  $57^{\mathrm m}$  &  $31^{\mathrm s}$  &  $4^{\mathrm u}$  &  $15^{\mathrm m}$  &  $39^{\mathrm s}$  \\
Bremen  &  $1$  &  $7$  &  $13$  &  $2$  &  $33$  &  $38$  &  $3$  &  $54$  &  $5$  \\
Breslau  &  $1$  &  $53$  &  $36$  &  $3$  &  $18$  &  $29$  &  $4$  &  $34$  &  $54$  \\
Cassel  &  $1$  &  $13$  &  $39$  &  $2$  &  $40$  &  $30$  &  $4$  &  $0$  &  $54$  \\
Cöln  &  $1$  &  $1$  &  $6$  &  $2$  &  $28$  &  $52$  &  $3$  &  $50$  &  $16$  \\
Danzig  &  $1$  &  $55$  &  $32$  &  $3$  &  $19$  &  $5$  &  $4$  &  $34$  &  $22$  \\
Göttingen  &  $1$  &  $15$  &  $51$  &  $2$  &  $42$  &  $28$  &  $4$  &  $2$  &  $39$  \\
Gotha  &  $1$  &  $20$  &  $59$  &  $2$  &  $47$  &  $34$  &  $4$  &  $7$  &  $30$  \\
Halle  &  $1$  &  $26$  &  $26$  &  $2$  &  $52$  &  $20$  &  $4$  &  $11$  &  $45$  \\
Hamburg  &  $1$  &  $12$  &  $26$  &  $2$  &  $38$  &  $19$  &  $3$  &  $58$  &  $14$  \\
\end{tabular}
\end{center}
\begin{center}\small\setlength{\tabcolsep}{3pt}
\begin{tabular}{l|rrr|rrr|rrr}
  &  \multicolumn{3}{c|}{Anfang.}  &  \multicolumn{3}{c|}{Mitte.}  &  \multicolumn{3}{c}{Ende.}  \\
\hline
 Hannover  &  $1^{\mathrm u}$  &  $13^{\mathrm m}$  &  $11^{\mathrm s}$  &  $2^{\mathrm u}$  &  $39^{\mathrm m}$  &  $36^{\mathrm s}$  &  $3^{\mathrm u}$  &  $59^{\mathrm m}$  &  $45^{\mathrm s}$ \\
 Königsberg  &  $2$  &  $4$  &  $43$  &  $3$  &  $27$  &  $20$  &  $4$  &  $41$  &  $32$ \\
 Leipzig  &  $1$  &  $28$  &  $38$  &  $2$  &  $54$  &  $25$  &  $4$  &  $13$  &  $45$ \\
 Mannheim  &  $1$  &  $11$  &  $47$  &  $2$  &  $39$  &  $36$  &  $4$  &  $0$  &  $31$ \\
 Marburg  &  $1$  &  $11$  &  $23$  &  $2$  &  $38$  &  $40$  &  $3$  &  $59$  &  $18$ \\
 München  &  $1$  &  $30$  &  $40$  &  $2$  &  $57$  &  $47$  &  $4$  &  $17$  &  $33$ \\
 Ofen  &  $2$  &  $11$  &  $2$  &  $3$  &  $35$  &  $49$  &  $4$  &  $51$  &  $31$ \\
 Prag  &  $1$  &  $41$  &  $51$  &  $3$  &  $7$  &  $50$  &  $4$  &  $25$  &  $46$ \\
 Tübingen  &  $1$  &  $16$  &  $42$  &  $2$  &  $44$  &  $36$  &  $4$  &  $5$  &  $24$ \\
 Wien  &  $1$  &  $55$  &  $45$  &  $3$  &  $21$  &  $24$  &  $4$  &  $38$  &  $28$ \\
\end{tabular}
\end{center}

Die Zeiten sind hier und in der folgenden Tabelle in \emph{wahrer Sonnenzeit} angegeben.

\begin{center}\small\setlength{\tabcolsep}{3pt}
\begin{tabular}{l|r|r|rrr|rrr}
 &  $P$  &  $D$  &  \multicolumn{3}{c|}{\shortstack{Erste innere\\Berührung.}}  &  \multicolumn{3}{c|}{\shortstack{Zweite innere\\Berührung.}}  \\
\hline
Berlin  &  $72^\circ$  &  $84''$  &  ---  &  ---  &   &  ---  &  ---  &   \\
Bremen  &  $65$  &  $9$  &  $2^{\mathrm u}$  &  $30'$  &  $46''$  &  $2^{\mathrm u}$  &  $36^{\mathrm m}$  &  $30^{\mathrm s}$  \\
Breslau  &  $78$  &  $131$  &  ---  &  ---  &   &  ---  &  ---  &   \\
Cassel  &  $66$  &  $1$  &  $2$  &  $37$  &  $37$  &  $2$  &  $43$  &  $24$  \\
Cöln  &  $62$  &  $55$  &  $2$  &  $27$  &  $23$  &  $2$  &  $30$  &  $21$  \\
Danzig  &  $79$  &  $186$  &  ---  &  ---  &   &  ---  &  ---  &   \\
Göttingen  &  $67$  &  $10$  &  $2$  &  $39$  &  $37$  &  $2$  &  $45$  &  $19$  \\
Gotha  &  $68$  &  $18$  &  $2$  &  $44$  &  $47$  &  $2$  &  $50$  &  $20$  \\
Halle  &  $70$  &  $48$  &  $2$  &  $50$  &  $20$  &  $2$  &  $54$  &  $19$  \\
Hamburg  &  $66$  &  $36$  &  $2$  &  $35$  &  $55$  &  $2$  &  $40$  &  $43$  \\
\end{tabular}
\end{center}
\begin{center}\small\setlength{\tabcolsep}{3pt}
\begin{tabular}{l|r|r|rrr|rrr}
  &  $P$  &  $D$  &  \multicolumn{3}{c|}{\shortstack{Erste innere\\Berührung.}}  &  \multicolumn{3}{c}{\shortstack{Zweite innere\\Berührung.}}  \\
\hline
 Hannover  &  $66^\circ$  &  $17''$  &  $2^{\mathrm u}$  &  $36^{\mathrm m}$  &  $49^{\mathrm s}$  &  $2^{\mathrm u}$  &  $42^{\mathrm m}$  &  $23^{\mathrm s}$ \\
 Königsberg  &  $81$  &  $217$  &  ---  &  ---  &   &  & --- &  \\
 Leipzig  &  $71$  &  $52$  &  $2$  &  $52$  &  $49$  &  $2$  &  $56$  &  $1$ \\
 Mannheim  &  $64$  &  $44$  &  $2$  &  $37$  &  $28$  &  $2$  &  $41$  &  $43$ \\
 Marburg  &  $65$  &  $18$  &  $2$  &  $35$  &  $53$  &  $2$  &  $41$  &  $27$ \\
 München  &  $71$  &  $2$  &  $2$  &  $54$  &  $54$  &  $3$  &  $0$  &  $40$ \\
 Ofen  &  $84$  &  $123$  &  ---  &  ---  &   &  ---  &  ---  &  \\
 Prag  &  $75$  &  $75$  &  ---  &  ---  &   &  ---  &  ---  &  \\
 Tübingen  &  $66$  &  $46$  &  $2$  &  $42$  &  $35$  &  $2$  &  $46$  &  $37$ \\
 Wien  &  $79$  &  $87$  &  ---  &  ---  &   &  ---  &  ---  &  \\
\end{tabular}
\end{center}

Hier bezeichnet $P$ den Punkt am Sonnenrand, an dem der Mond zu Beginn der Finsternis zuerst eingreift, wobei man die Grade vom nördlichsten Punkt (bezogen auf den Vertikalkreis) nach \SourceLeaf{629} Westen zählt. $D$ hingegen gibt in Bogensekunden den scheinbaren Abstand zwischen den Mittelpunkten von Sonne und Mond zur Mitte der Finsternis an. Bei allen Rechnungen wird eine Abplattung von $\tfrac{1}{300}$ angenommen.

Die wichtigsten Zeitpunkte der Finsternis für andere Orte in Deutschland lassen sich mit den folgenden Formeln näherungsweise aus den Zeitpunkten ableiten, die in der ersten Tabelle oben für einen nahe gelegenen Ort angegeben sind. Dabei bezeichnet d$L$ den Längenunterschied zwischen beiden Orten in Zeitsekunden, d$B$ den Breitenunterschied in Bogensekunden und d$Z$ die Änderung, die zu den Tabellenwerten addiert werden muss. d$L$ ist positiv anzusetzen, wenn der Ort, für den man die Zeitpunkte bestimmen möchte, östlich des aus der Tabelle gewählten Ortes liegt, und d$B$ ist ebenfalls positiv anzusetzen, wenn jener nördlich von diesem liegt.

Für den Anfang
\[
\mathrm{d}Z = +1.313\,\mathrm{d}L - 0.0296\,\mathrm{d}B
\]
Für die Mitte
\[
\mathrm{d}Z = +1.242\,\mathrm{d}L - 0.0343\,\mathrm{d}B
\]
Für das Ende
\[
\mathrm{d}Z = +1.118\,\mathrm{d}L - 0.0363\,\mathrm{d}B
\]

Besonders wichtig ist es, die Grenzlinien der Zone zu kennen, innerhalb derer die Finsternis ringförmig erscheint. Herr Prof. \textsc{Gerling} hat die folgenden Punkte auf diesen Grenzlinien durch Rechnung bestimmt, wobei die Längengrade von Ferro aus gerechnet sind.

\begin{center}
\begin{tabular}{c|c|c|c|c}
\shortstack{Wahre\\Zeit in\\Paris.} & \multicolumn{2}{c|}{\shortstack{Berührung der\\südl. Ränder}} & \multicolumn{2}{c}{\shortstack{Berührung der\\nördl. Ränder.}} \\
 & Breite. & Länge. & Breite. & Länge. \\
\hline
$2^{\mathrm u}\ 4^{\mathrm m}$ & & & $52^\circ\ 59'$ & $22^\circ\ 43'$ \\
6 & & & $52\ \ 15$ & $23\ \ 15$ \\
8 & $54^\circ\ 12'$ & $28^\circ\ 50'$ & $51\ \ 31$ & $23\ \ 46$ \\
10 & $53\ \ 26$ & $29\ \ 18$ & $50\ \ 47$ & $24\ \ 17$ \\
12 & $52\ \ 40$ & $29\ \ 47$ & $50\ \ \ 4$ & $24\ \ 48$ \\
14 & $51\ \ 55$ & $30\ \ 16$ & $49\ \ 21$ & $25\ \ 19$ \\
16 & $51\ \ 10$ & $30\ \ 46$ & $48\ \ 39$ & $25\ \ 49$ \\
18 & $50\ \ 26$ & $31\ \ 15$ & $47\ \ 57$ & $26\ \ 19$ \\
20 & $49\ \ 42$ & $31\ \ 45$ & $47\ \ 15$ & $26\ \ 49$ \\
22 & $48\ \ 58$ & $32\ \ 15$ & $46\ \ 34$ & $27\ \ 20$ \\
24 & $48\ \ 15$ & $32\ \ 46$ & $45\ \ 53$ & $27\ \ 50$ \\
26 & $47\ \ 32$ & $33\ \ 17$ & & \\
28 & $46\ \ 50$ & $33\ \ 49$ & & \\
\end{tabular}
\end{center}

Verbindet man diese Punkte auf einer Landkarte durch gerade Linien, so trifft man auf die folgenden Orte, von denen keiner mehr als etwa $2\tfrac{1}{2}$ Meilen von den Grenzlinien entfernt liegt. Die typografisch hervorgehobenen Orte liegen der Grenze am nächsten.

Für die Berührung der \emph{südlichen} Ränder: Heiligenhafen, Neustadt, Grevesmühlen, Wismar, Schwerin, Parchim, Grabow, \emph{Perleberg}, Werben, Havelberg, \emph{Arneburg}, Tangermünde, Rathenow, Genthin, Ziesar, Coswig, Wittenberg, \emph{Kemberg}, Düben, Torgau, Wurzen, Oschatz, Döbeln, \emph{Freiberg}, Seyda, Kommotau, \emph{Rakonitz}, Przibram, \emph{Teyn}, Budweis, \emph{Grazen}, \emph{Ips}, \emph{Oberndorf}, Tirnitz, Fürstenfeld u.\,s.\,w.

Für die Berührung der \emph{nördlichen} Ränder: Staveren, Medenblick, Harderwyk, Amersfort, \emph{Barneveld}, Arnheim, \emph{Nimwegen}, Cleve, Grave, \emph{Genep}, Geldern, \emph{Venloo}, Crefeld, Roermunde, \emph{Erkelens}, Jülich, \emph{Nidegen}, Münster-Eiffel, Aldenau, Mondheim, \emph{Zell}, Trarbach, Kirchberg, \emph{Oberstein}, Münchweiler, \SourceLeaf{630} Kaiserslautern, Bergzabern, Hagenau, \emph{Stollhofen}, Achern, Freydenstadt, Hornberg, Rottweil, Villingen, \emph{Geisingen}, Schaffhausen, \emph{Stein}, \emph{Frauenfeld}, Winterthur, Chur u.\,s.\,w.

Es ist sehr zu wünschen, dass in der Nähe der Grenzlinien zahlreiche Beobachtungen angestellt werden. Sie sind schon dann von Wert, wenn der Beobachtungsort angegeben und vermerkt wird, ob sich ein Ring gebildet hat oder nicht und, falls ja, wie lange er zu sehen war. Da es hier nur um die Messung einer sehr kurzen Zeitspanne geht, kann man sich beim Fehlen einer Sekundenuhr notfalls damit behelfen, die Schläge einer gewöhnlichen Taschenuhr oder die Schwingungen einer schweren, an einem Faden aufgehängten Kugel zu zählen. Voraussetzung ist, dass man im ersten Fall die Zahl der Schläge pro Minute und im zweiten die Länge des Fadens und die Größe der Kugel in einem bekannten Maß mit angibt. Eine solche Beobachtung müsste jedoch mit einem guten Fernrohr angestellt werden, wenn sie von Wert sein soll.

Da selbst an den Orten, an denen die Finsternis ringförmig erscheint, noch ein volles Achtel der Sonnenscheibe unverfinstert bleibt, ist eine starke Abnahme des Tageslichts nicht zu erwarten, noch weniger, dass Sterne mit bloßem Auge sichtbar werden. Für diejenigen, die während der größten Verfinsterung versuchen möchten, einige der hellsten und der Sonne nächsten Fixsterne mit Fernrohren aufzusuchen oder mit guten Reflexionsinstrumenten ihre Abstände vom Mondrand zu messen, hat Herr Prof. \textsc{Gerling} noch die folgenden Angaben beigefügt, die das Auffinden dieser Sterne erleichtern können.

\begin{center}
\begin{tabular}{lcc}
Für Regulus & $P = 65\tfrac{3}{4}^\circ$ westl., & $D = 17\tfrac{1}{2}^\circ$ \\
Für $\ObservedBeta$ Löwe & $P = 41\tfrac{3}{4}^\circ$ östl., & $D = 13^\circ$ \\
Für Spica & $P = 116^\circ$ östl., & $D = 36\tfrac{1}{2}^\circ$. \\
\end{tabular}
\end{center}

Hier bezeichnet $D$ die Entfernung des Sterns vom Mittelpunkt der Sonne und $P$ den Winkel, den der Großkreis vom Sonnenmittelpunkt zum Stern mit dem Großkreis vom Sonnenmittelpunkt zum Nordpol bildet. Die Angaben sind für $2^{\mathrm u}\ 15^{\mathrm m}$ W.\,Z. in Paris berechnet, können aber in ganz Deutschland für die Zeit der größten Verfinsterung verwendet werden.
\par\clearpage
